\documentclass[10pt,a4paper]{amsart}
\usepackage[margin=19mm]{geometry}
\usepackage{amsmath,amssymb,mathtools}
\usepackage[scr=rsfs,cal=euler]{mathalfa}
\usepackage{xfrac}
\usepackage[unicode,hidelinks,bookmarks=false]{hyperref}
\newcommand{\ie}{i.e.\ }
\newcommand{\cf}{cf.\ }
\newcommand{\resp}{respectively\ }
\newcommand{\Subsection}{\subsection}
\providecommand{\nobreakdash}{\nobreak\hbox{-}}
\setlength{\emergencystretch}{2em}
\multlinegap0pt
\usepackage{amssymb}
\newtheorem{theorem}{Theorem}
\title{{\bf Critical Sets in Latin Squares\\ and\\ Associated Structures}}
\author{Richard Winston Bean, B.Sc. (Hons) (UQ)}
\date{Source: arXiv:2606.12996v2 (CC BY 4.0)}
\begin{document}
\maketitle

\noindent\emph{Source and licence.} {\bf Critical Sets in Latin Squares\\ and\\ Associated Structures}. Version arXiv:2606.12996v2 (CC BY 4.0), distributed by arXiv under the Creative Commons Attribution 4.0 International licence, http://creativecommons.org/licenses/by/4.0/, as recorded in the arXiv licence metadata for this exact version and shown on the abstract page https://arxiv.org/abs/2606.12996v2. Full licence notice: \texttt{licenses/arxiv-2606.12996-CC-BY-4.0.txt}.\par\smallskip

\noindent\textit{Editorial scope.} Counted unit: the article's theorem theorem (source0.tex lines 3008--3011). The article's complete original proof of it is delivered with it (source0.tex lines 3013--3044, one \texttt{proof} environment of 32 lines). No mathematics is written, repaired or re-derived: the statement and the proof are the article's own lines, in the article's own notation, and no retained line is substituted or re-flowed.  No further source-local context is needed: every cross-reference of every retained line resolves inside the two delivered spans.  Nothing was omitted from any retained span, so the package carries no omission record.  No retained line contains a citation, so the wrapper carries no bibliography and no bibliography entry.  No retained line contains a figure environment, an image-inclusion command or drawing code, so no image is delivered and the package declares no asset dependency.  Editorial formatting only, in addition: an AMS \texttt{amsart} preamble that restates the article's own declarations and macros used by the retained lines, namely the environments \texttt{theorem} on separate counters, together with the standard packages those lines need.  The retained environments print in this document as Theorem 1 in order of appearance, whereas the article prints the same items with the numbering of its own full document.  The complete native source of the article as distributed in the arXiv e-print for this exact version is bundled unchanged as \texttt{sources/202666/source0.tex}.
\begin{theorem}
For $\alpha \geq 2$,
$I(2^\alpha m) \geq (2^\alpha m)^2 (3m 2^\alpha+2^\alpha-4)/16$.
\end{theorem}

\begin{proof}
{
Consider ${\mathbb Z}_2^2$ displayed above; it contains 12 distinct intercalates.
Then if $\{(r_1,c_1;e_1),(r_1,c_2;e_2),(r_2,c_1;e_2),(r_2,c_2;e_1)\}$
is an intercalate in ${\mathbb Z}_2^2$, we have that
\begin{eqnarray*}
D &=& \{ (i,j;L'_{ij}) \mid ((c_1m \leq j \leq c_1m+m-1) \vee (c_2m \leq j \leq c_2m+m-1)) \wedge \\
&& ((r_1m \leq i \leq r_1m+m-1) \vee (r_2m \leq i \leq r_2m+m-1)) \}
\end{eqnarray*}
is a subsquare of $L'$ which is isomorphic to one of $L_1, L_2, L_3, L_4, L_5$
or $L_6$ and thus contains exactly $m^3$ intercalates.  Therefore $I(L') = 12m^3$, and thus $I(4m) \geq 12m^3$.

Heinrich and Wallis counted the number of intercalates in the direct
product of two Latin squares $M$ (of order $k$) and $N$ (of order $l$).
This count was used to create a new lower bound on $I(kl)$:
\begin{eqnarray*}
I(kl) \geq I(k) l^2 + I(l) k^2 + 4.I(k).I(l)
\end{eqnarray*}

If we use the Latin squares $M = {\mathbb Z}_2^{\alpha-2}$ and $N = L'$,
we may use this formula with $k = 2^{\alpha-2}$ and $l = 4m$
and the values $I(2^{\alpha-2}) = \displaystyle{\frac{(2^{\alpha-2})^2
(2^{\alpha-2}-1)}{4}}$ (known from Heinrich and Wallis), and $I(4m)
\geq 12m^3$.  Thus we deduce that:
\begin{eqnarray*}
I(2^\alpha m) &\geq& I(4m).(2^{\alpha-2})^2 + I(2^{\alpha-2}).(4m)^2 + 4.I(2^{\alpha-2}).I(4m) \\
&=& 12m^3.2^{2\alpha-4} + 16m^2.\displaystyle{\frac{(2^{\alpha-2})^2 (2^{\alpha-2}-1)}{4}} +  \\
&& 4.\displaystyle{\frac{(2^{\alpha-2})^2 (2^{\alpha-2}-1)}{4}}.12m^3 \\
&=& (2^\alpha m)^2 (3m.2^\alpha + 2^\alpha - 4)/16
\end{eqnarray*}
}
\end{proof}

\end{document}
